\section{Implementation Status}

ActPlane is implemented as a Rust collector and an eBPF loader/program pair.
The current codebase builds a single \texttt{actplane} binary that embeds the
compiled loader, extracts it at runtime, compiles project policy YAML, and runs
a target command under the policy harness.

\subsection{Runner}

The command \texttt{actplane run -- <cmd>} discovers \texttt{actplane.yaml}
upward from the current directory unless \texttt{--policy} or \texttt{--rule} is
specified. The policy YAML contains a \texttt{policy: |} DSL block and an
optional feedback path. Run mode requires an \texttt{AGENT} label; before the
target begins executing, the loader seeds the stopped target pid with that label
so all descendants inherit the agent identity.

The runner prepares \texttt{.actplane/last-violation.txt} and a hook state file,
exports their paths to the child, waits until the loader reports readiness, then
continues the target process. This ensures the first target action is observed.

\subsection{BPF Backends}

The eBPF program attaches to process lifecycle tracepoints for fork, exec, and
exit, to syscall tracepoints for observation-compatible file and connect events,
and to BPF-LSM hooks for pre-operation enforcement. The loader checks whether
the running kernel has the \texttt{bpf} LSM active. If not, it uses tracepoint
mode. In tracepoint mode, supported effects are exactly \texttt{audit} and
\texttt{kill}; \texttt{block} is unsupported because there is no pre-operation
verdict to return.

\subsection{Current Limitations}

The implementation deliberately exposes several limitations rather than hiding
them behind ambiguous semantics.

\textbf{Path identity.} File labels are currently keyed by path hash. This is
sufficient for the prototype but should move toward inode/device identity for a
stronger claim.

\textbf{Endpoint identity.} Kernel connect matching is numeric IPv4. Hostname,
DNS, SNI, and HTTP-layer policies require userspace assist or additional
instrumentation.

\textbf{Exec arguments.} \texttt{@arg} predicates are currently observed after
exec unless an argv cache is added before \texttt{bprm\_check\_security}. Thus
\texttt{effect block} with an argv predicate is not available through the
tracepoint path; \texttt{effect audit} can report and \texttt{effect kill} can
terminate after exec.

\textbf{Precision.} Labels propagate at object granularity. This can over-taint
compared with byte-level DIFT systems. The evaluation must therefore measure
false positives and validate declassification/gate paths, not only violation
coverage.

\subsection{Engineering Invariants}

The codebase now follows three invariants that matter for the paper:

\begin{itemize}
\item The only kernel output for policy violations is
\texttt{TAINT\_VIOLATION}; ordinary tracing noise is not part of the harness
interface.
\item Each rule explicitly declares its effect. There is no global fallback that
converts \texttt{block} into \texttt{kill}.
\item The post-tool feedback adapter is a transport bridge, not a second policy
engine.
\end{itemize}
